\chapter*{Abstract}
\addcontentsline{toc}{chapter}{Abstract}

Stress-related strain is increasingly common, yet many people find it difficult to sustain meditation practice because conventional audio-guided formats can feel repetitive, provide limited personalization, and offer little support when attention drifts. Virtual reality (VR) has the potential to reduce external distractions through immersive environments, while AI-driven voice interaction and adaptive audiovisual feedback may deliver more engaging and responsive guidance without requiring a human instructor.

This thesis investigates whether an AI-guided, voice-interactive VR meditation system can enhance relaxation and emotional well-being through adaptive audiovisual cues. A meditative application was developed in Unity for the HTC VIVE Pro headset. The system provides meditation guidance via GPT-generated voice prompts within a calming virtual environment and incorporates physiological monitoring to support responsive transitions. To evaluate the approach, a controlled study with 44 participants compared VR-based AI-guided meditation with a conventional meditation condition. Physiological indicators of arousal and relaxation, including heart rate variability and skin conductance, were recorded using the Biofeedback 2000 system and complemented by self-report measures of relaxation and emotional state.

Overall, the results indicate that AI-guided VR meditation enriched with adaptive audiovisual cues can promote stronger relaxation and more positive emotional experiences than conventional methods. The findings underscore the potential of adaptive, multi-modal VR meditation systems as supportive tools for stress reduction and mental well-being, and they motivate further research into personalization strategies and long-term effects in real-world use.

\vspace{1cm}
\textbf{Keywords:} VR, Virtual Reality, Meditation, Artificial Intelligence, Biofeedback, Relaxation, Emotional well being
\clearpage
\chapter*{Zusammenfassung}
\addcontentsline{toc}{chapter}{Zusammenfassung}

Stressbedingte Belastungen nehmen zunehmend zu, dennoch fällt es vielen Menschen schwer, eine regelmäßige Meditationspraxis aufrechtzuerhalten. Konventionelle, audioangeleitete Formate werden häufig als repetitiv erlebt, bieten nur begrenzte Personalisierung und unterstützen wenig, wenn die Aufmerksamkeit abschweift. Virtuelle Realität (VR) kann durch immersive Umgebungen externe Ablenkungen reduzieren, während KI-gestützte Sprachinteraktion und adaptive audiovisuelle Rückmeldungen eine ansprechendere und responsivere Anleitung ermöglichen können – ohne die kontinuierliche Präsenz einer menschlichen Lehrperson.

Diese Arbeit untersucht, ob ein KI-geleitetes, sprachinteraktives VR-Meditationssystem durch adaptive audiovisuelle Hinweise die Entspannung und das emotionale Wohlbefinden verbessern kann. Hierzu wurde eine Meditationsanwendung in Unity für das HTC VIVE Pro Headset entwickelt. Das System bietet Meditationsanleitungen mittels GPT-generierter Sprachprompts innerhalb einer beruhigenden virtuellen Umgebung und integriert physiologisches Monitoring, um responsive Übergänge zu unterstützen. Zur Evaluation wurde eine kontrollierte Studie mit 44 Teilnehmenden durchgeführt, in der eine VR-basierte, KI-geleitete Meditation mit einer konventionellen Meditationsbedingung verglichen wurde. Physiologische Indikatoren für Aktivierung und Entspannung, darunter die Herzratenvariabilität und die Hautleitfähigkeit, wurden mit dem Biofeedback-2000-System erfasst und durch Selbstberichte zu Entspannung und emotionalem Zustand ergänzt.

Insgesamt deuten die Ergebnisse darauf hin, dass KI-geleitete VR-Meditation, angereichert durch adaptive audiovisuelle Hinweise, eine stärkere Entspannung und positivere emotionale Erfahrungen fördern kann als konventionelle Methoden. Die Befunde unterstreichen das Potenzial adaptiver, multimodaler VR-Meditationssysteme als unterstützende Werkzeuge zur Stressreduktion und Förderung des mentalen Wohlbefindens und motivieren weitere Forschung zu Personalisierungsstrategien sowie langfristigen Effekten im realen Anwendungskontext.

\vspace{1cm}
\textbf{Schlüsselwörter:} VR, Virtual Reality, Meditation, Künstliche Intelligenz, Biofeedback, Entspannung, Emotionales Wohlbefinden
